% Source: https://artofproblemsolving.com/community/c3h2852425_some_problems__by_megahertz13
% Stable identity: aops-topic-2852425-q03
Rotate $\triangle APC$ by $60^\circ$ about $A$ so that $C$ goes to $B$; let $P$ go to $P'$. Then $AP'=AP=1$ and $\angle PAP'=60^\circ$, so $\triangle APP'$ is equilateral with $PP'=1$; and $BP'=CP=2$.
\begin{asy}
usepackage("amsmath");
size(4.5cm);
defaultpen(fontsize(10pt));
real s=sqrt(7);pair A=(0,0),B=(s,0),C=s*dir(60);
real ax=(1-3+7)/(2*s);pair P=(ax,sqrt(1-ax^2));
pair P2=A+rotate(-60)*(P-A);
draw(A--B--C--cycle);draw(P--A);draw(P--B);draw(P--C);draw(P--P2--A,dashed);draw(P2--B,dashed);
dot(P);label("$A$",A,SW);label("$B$",B,SE);label("$C$",C,N);label("$P$",P,N);label("$P'$",P2,S);
\end{asy}

In $\triangle PP'B$: $PP'^2+PB^2=1+3=4=P'B^2$, so $\angle P'PB=90^\circ$. Therefore $\angle APB=60^\circ+90^\circ=150^\circ$, and by the law of cosines
\[s^2=AB^2=1+3-2\sqrt3\cos150^\circ=4+3=7,\qquad s=\boxed{\sqrt7}.\]
